% Separate post-CP25 candidate. Exact point-field proofs precede finite checks.
\subsection{Marked low lines and a larger exact tame sector}

Throughout the local lemmas \(p\ge5\), \(q=p^2\), \(N=q-1\),
and \(K\) is complete discretely valued of mixed characteristic
\((0,p)\), with algebraically closed residue field, \(\mu_p\subset K\),
and absolute ramification \(e\) prime to \(p\). Put \(s=e(p+1)\).
The actual good field \(L/K\) is tame of degree \(N\), with
\(v_L(p)=eN=(p-1)s\), and \(s\) is a multiple of \(N\).
For a two-slope group use the notation of B69:
\[
 0<B<pe,\quad B\equiv c\pmod{p+1},\quad
 \alpha=s-B,\quad\beta=B/p,\quad
 \lambda_H=\zeta^{c-B}.
 \tag{B-RD71.1}
\]
All coordinate valuations below belong to actual integral good
models and their finite-flat quotients.

\begin{lemma}[Prime Newton parameters are determined at the first level]
\label{lem:B-RD71-first-level}
\status{Candidate; complete proof supplied, independent audits pending.}
Two two-slope groups with parameters \(B_1,B_2\) prime to \(p\)
and isomorphic actual \(p\)-torsion have \(B_1=B_2\) and the
same finite special inertia character.
\end{lemma}
\begin{proof}
The common first-level division field \(R/L\) is cyclic of
degree \(p\). Indeed the high character and determinant are
trivial over \(L\), so the residual action lies in the
upper-unipotent group of order \(p\); the low valuation \(B_i/p\)
is not integral over \(L\), so the action is nontrivial.
The high lines are the common unique wild fixed line.
In \(R\), a low point has coordinate valuation \(B_i\).
For any nonidentity \(g\in\operatorname{Gal}(R/L)\), its
displacement is a nonzero high point, of valuation
\(p(s-B_i)\). Formal subtraction has the same valuation as
coordinate subtraction, since the relevant parameters have
positive valuation.
The exact prime-valuation part of
Lemma~\ref{lem:B-RD60-break} therefore gives the actual lower break
\[
 r=p(s-B_i)-B_i=ps-(p+1)B_i.
 \tag{B-RD71.2}
\]
The field and break are common, so \(B_1=B_2\).
Their identified high-line characters in (B-RD71.1) then
give \(c_1=c_2\pmod N\). This argument uses only the first-level
torsion isomorphism; the second-level assumption in the broader
comparison lemma is unnecessary for this particular branch.
\end{proof}

\begin{lemma}[The preserved low line in a small high quotient]
\label{lem:B-RD71-marked-low}
\status{Candidate; complete proof supplied, independent audits pending.}
Let \(M\) have two-slope parameter \(B<e\) and high line \(H\).
The actual quotient \(M'=M/H\) has parameter \(D=pB\), and
\[
 J=\operatorname{im}\bigl(M[p]\longrightarrow M'[p]\bigr)
   \simeq M[p]/H
 \tag{B-RD71.3}
\]
is an inertia-stable low line of \(M'[p]\), of valuation \(B\).
The first-level module \(M'[p]\) is rational over \(L\).
An actual \(p^3\)-torsion isomorphism identifying two original
high lines induces an actual \(p^2\)-torsion isomorphism of
the quotients which also identifies these marked low lines.
The quotient's low/high eigenvalue ratio is
\(\zeta^{(p+1)B}\); it may be scalar.
\end{lemma}
\begin{proof}
The high-quotient product calculation in
Lemma~\ref{lem:B-RD69-quotients} gives \(D=pB\).
The \(p-1\) nonzero points of \(M[p]/H\) have quotient
valuation \(B\), whereas the high valuation in \(M'[p]\)
is \(s-pB>B\), since \(B<e\).
Thus their image is precisely a low line, and is stable
under the actual generic inertia action.
The high line of \(M'[p]\) is stable too, because it is
characterized by its strictly larger valuation.
These two distinct invariant \(\mathbf F_p\)-lines make
the entire residual representation diagonal with
\(\mathbf F_p^\times\)-valued characters. They are killed
by the tame extension of degree \(N\); consequently \(M'[p]\)
is \(L\)-rational, even when the original residual action is wild.

The actual quotient has tangent exponent \(pc\).
Its high eigenvalue is \(\zeta^{pc-pB}\), and its other
eigenvalue is \(\zeta^{c+pB}\), by its determinant and
\(p^2\equiv1\pmod N\). Their ratio has exponent
\((1-p)c+2pB\), congruent to \((p+1)B\), because
\((p-1)(c-B)\) is divisible by \(N\).
In particular scalar action can occur when \(p-1\mid B\).
No uniqueness of an unmarked complementary line is assumed.

The induced actual \(p^2\)-isomorphism is obtained from
\(\{P\in M[p^3]:p^2P\in H\}/H\).
It carries the image of the intrinsic subgroup \(M[p]\)
to the corresponding image in the other quotient. This
proves preservation of \(J\) as an actual Galois-module
subgroup. The finite-flat quotients and semilinear generic
descent are those of Lemma~\ref{lem:B-RD69-quotients}.
\end{proof}

\begin{proposition}[A scalar-safe recovery of the quotient high line]
\label{prop:B-RD71-separated-high}
\status{Candidate; complete proof supplied, independent audits pending.}
Suppose an actual \(p^3\)-torsion isomorphism identifies the
high lines of two two-slope groups, both with
\[
 B_i=pb_i<e.
 \tag{B-RD71.4}
\]
If \(q\min(b_1,b_2)>e\), then \(B_1=B_2\) and their special
characters agree. This assertion includes scalar residual
actions in the first quotients.
\end{proposition}
\begin{proof}
If \(b_1=b_2\), the common high-line character already
identifies \(c_1=c_2\). Otherwise interchange the models
so that \(b_1<b_2\).
Quotient by the common high line and write \(M_i'=M_i/H_i\).
By Lemma~\ref{lem:B-RD71-marked-low} these groups have an
actual \(p^2\)-torsion isomorphism \(\phi'\), rational
first-level modules over \(L\), and parameters
\[
 D_i=q b_i,\qquad D_i>e,\qquad
 A_i=s-D_i,\qquad C_i=D_i/p .
 \tag{B-RD71.5}
\]
Here \(A_i\) and \(C_i\) are their high and low first-level
valuations. In particular \(p\nmid A_i\).

We prove that \(\phi'\) identifies their high lines.
Suppose otherwise. Choose a primitive order-\(p^2\) point
\(X_1\) of \(M_1'\) whose \(p\)-multiple is a nonzero high
point, and put \(X_2=\phi'(X_1)\).
The \(p\)-multiple of \(X_2\) is then low in \(M_2'\).
Because \(D_1>e\), the high-target inverse Newton polygon
of \(M_1'\) has a single segment of degree \(q\):
the exact inverse calculation in Lemma~\ref{lem:B-RD69-quotients}
gives
\[
 v_L(X_1)=A_1/q,\qquad
 v_L(X_2)=D_2/p^3=C_2/q .
 \tag{B-RD71.6}
\]
The second identity is the low-target inverse calculation.
Since \(A_1=s-qb_1\) is prime to \(p\), the denominator of
the first valuation is exactly \(q\).

Both \(p\)-multiples are \(L\)-rational, and both first-level
modules are \(L\)-rational. Therefore each point field over
\(L\) is Galois with group a subgroup of its first-level
translation group, of order at most \(q\): every conjugate
is \(X_i+Q\), and all such sums lie in the point field.
The actual torsion isomorphism identifies their stabilizers,
so the two point fields are the same field \(E\).
The first valuation denominator forces its ramification
index to be at least \(q\). Consequently
\[
 [E:L]=e(E/L)=q,\qquad
 v_E(X_1)=A_1,\quad v_E(X_2)=C_2,
 \tag{B-RD71.7}
\]
and every first-level translation occurs.

Choose \(g\in\operatorname{Gal}(E/L)\) which translates \(X_1\)
by a nonzero high point of \(M_1'[p]\).
Its corresponding translation in \(M_2'[p]\) is low, by
the assumed failure to identify high lines.
Thus its coordinate displacements have valuations
\(qA_1\) and \(qC_2\), respectively.
Since \(p\nmid A_1\), the exact lower-break lemma gives
\(b(g)=(q-1)A_1\).
The unconditional lower bound on the second coordinate
then requires
\[
 qC_2\ge C_2+(q-1)A_1,\qquad\text{hence }C_2\ge A_1.
 \tag{B-RD71.8}
\]
Yet \(pb_2=B_2<e\), so
\[
 \begin{aligned}
 A_1-C_2
   &=s-qb_1-pb_2\\
   &=s-p(p+1)b_2+q(b_2-b_1)
     >q(b_2-b_1)>0.
 \end{aligned}
 \tag{B-RD71.9}
\]
This contradiction identifies the high lines, without
assuming distinct residual eigenvalues.

We may therefore quotient both \(M_i'\) by their high lines.
Since \(D_i>e\), their actual high quotients have parameters
\(s-D_i=A_i\), both prime to \(p\).
The induced comparison is now exactly at the first level:
\[
 (M_i'/H_i')[p]
   =\{Z\in M_i'[p^2]:pZ\in H_i'\}/H_i'.
 \tag{B-RD71.10}
\]
Lemma~\ref{lem:B-RD71-first-level} gives \(A_1=A_2\),
so \(D_1=D_2\), contrary to \(b_1<b_2\).
Thus only \(b_1=b_2\) was possible, proving the assertion.
\end{proof}

\begin{proposition}[Common-high comparisons in the enlarged sector]
\label{prop:B-RD71-common-high}
\status{Candidate; complete proof supplied, independent audits pending.}
Suppose now that
\[
 e=t(p-1),\qquad 2\le t\le2p+1,\qquad p\nmid t.
 \tag{B-RD71.11}
\]
An actual \(p^3\)-torsion isomorphism between two two-slope
groups which identifies their high lines determines their
finite special inertia characters.
\end{proposition}
\begin{proof}
Their common finite determinant and high-line character give
\[
 c_1\equiv c_2\pmod{p-1},\qquad
 c_1-B_1\equiv c_2-B_2\pmod N.
 \tag{B-RD71.12}
\]
Any case with a prime-to-\(p\) parameter is already settled
by Lemma~\ref{lem:B-RD69-rules}; it includes \(B_i=e\).
Assume both \(B_i\) are divisible by \(p\).
Neither equals \(e\). If both are above \(e\), their actual
high quotients have prime parameters \(s-B_i\), and the
inherited \(p^2\)-comparison recovers their characters.
If the parameters lie on opposite sides of \(e\), those
quotients have opposite \(p\)-divisibility and the same
comparison lemma excludes them.

It remains that \(B_i=pb_i<e\). The range (B-RD71.11) gives
\[
 e\le2p^2-p-1<2q,\qquad 1\le b_i\le2p-2.
 \tag{B-RD71.13}
\]
If the \(b_i\)'s are equal, (B-RD71.12) identifies \(c_1=c_2\).
Otherwise take \(b_1<b_2\).
If \(b_1\ge2\), then \(qb_1\ge2q>e\);
Proposition~\ref{prop:B-RD71-separated-high} forbids the
inequality \(b_1<b_2\).
Thus \(b_1=1\).
Equation (B-RD71.12) implies \(b_2\equiv1\pmod{p-1}\),
and the upper bound \(b_2\le2p-2\) forces \(b_2=p\).
Exactly as in (B-RD70.5),
\[
 c_2-c_1\equiv p(p-1),\qquad
 c_1\equiv p\pmod{p+1}
 \quad\Longrightarrow\quad c_2\equiv pc_1\pmod N.
 \tag{B-RD71.14}
\]
This exchanges the two Frobenius-conjugate eigenvalues
and preserves the special character. Every case is covered.
\end{proof}

\begin{proposition}[Tame residual comparisons in the enlarged sector]
\label{prop:B-RD71-tame}
\status{Candidate; complete proof supplied, independent audits pending.}
Under (B-RD71.11), two two-slope groups with tame residual
inertia and isomorphic actual \(p^3\)-torsion have equal
finite special inertia characters.
\end{proposition}
\begin{proof}
Write \(B_i=pb_i\), \(0<b_i<e\), since their first-level
modules are rational over \(L\).
If both \(b_i\) are prime to \(p\), the second rule of
Lemma~\ref{lem:B-RD69-rules} gives the conclusion.
If exactly one is prime, say \(b_1=b\), \(b_2=pk\), take
primitive points whose \(p\)-multiples are outside both
high lines. Their common point field over \(L\) has degree
\(q\), forced by the first valuation \(b/q\).
Their normalized coordinate valuations are \(b\) and \(pk\).
Distinct high lines would permit a translation high in
the first model and low in the second. The exact break
from the first coordinate and the lower bound for the second
would give
\[
 q(\alpha_1-pk)\le b-pk .
 \tag{B-RD71.15}
\]
But \(\alpha_1=s-pb\ge e+p\), \(pk<e\), and
\(e<2q<p^3\). The left side is greater than \(p^3\),
the right side less than \(p^3\). Thus their high lines
agree, and Proposition~\ref{prop:B-RD71-common-high} applies.

Finally let \(b_i=pk_i\), so \(B_i=qk_i\).
If the high lines agree, apply the same proposition.
If not, and one \(B_i>e\), quotient that model by its high
line and the other by the corresponding stable low line.
The inherited actual \(p^2\)-comparison has parameters
\(s-B_i\), prime to \(p\), and \(B_j/p=pk_j\),
divisible by \(p\). This is excluded by
Lemma~\ref{lem:B-RD69-rules}.
The equality \(B_i=e\) is impossible because \(p\nmid e\).
If both \(B_i<e<2q\), then \(B_1=B_2=q\).

For completeness the exchanged-line obstruction does not
depend on the particular index used in B70.
The residual high and low eigenvalue exponents in this
last case are \(c_i-q\) and \(pc_i+q\).
Semisimplicity of tame inertia and noncorresponding high
lines imply
\[
 c_1-q\equiv pc_2+q,\qquad
 c_2-q\equiv pc_1+q\pmod N.
 \tag{B-RD71.16}
\]
These equalities remain valid for scalar action as well.
Multiplying the first by \(p\) and adding the second would
give \(p-1\mid2q\), or \(p-1\mid2\), impossible for \(p\ge5\).
Hence the high lines must correspond after all.
This finishes the last case.
\end{proof}

\begin{theorem}[A larger exact tame sector in full height-two monodromy]
\label{thm:B-RD71}
\status{Candidate; complete proof supplied, independent audits pending.}
For every prime \(p\ge5\), in the full B.4 monodromy setting
with height two, perfect residue field and \(\mu_p\subset K\),
\[
 \boxed{\displaystyle
 \nu_{\min}\bigl(p,t(p-1),2\bigr)=2
 \quad(2\le t\le2p+1,\ p\nmid t).}
 \tag{B-RD71.17}
\]
For \(p=5\), the additionally settled prime-to-\(5\) indices
are \(e=32,36,44\).
\end{theorem}
\begin{proof}
After completed maximal unramified base change, compare
actual \(p^3\)-torsion. Since \(p\nmid e\), its cyclotomic
image has order \(p^2\); the determinant separates the
total dimensions \(0,1,2\) and their finite determinant
characters exactly as in Theorem~\ref{thm:B-RD69}.
Dimensions zero and two are detected residually and by
Cartier duality. Two systems filtered by potentially good
height-one pieces are also detected residually.
Mixed whole-supersingular/filtered comparisons are excluded
by Proposition~\ref{prop:B-RD63-mixed}, since \(e<2q<p^3\).

For two whole potentially supersingular groups,
Proposition~\ref{prop:B-RD69-one-slope} excludes mixed
Newton polygons and recovers two one-slope characters,
without an upper restriction on \(e\).
For two unequal polygons, residual wild inertia makes
their high lines the common unique fixed line, so
Proposition~\ref{prop:B-RD71-common-high} applies.
Residual tame inertia is covered by
Proposition~\ref{prop:B-RD71-tame}.
Thus actual \(p^3\)-torsion always determines the character.
For every index in (B-RD71.11), \(e\ge2(p-1)>p\).
Theorem~\ref{thm:B-RD63-arc}, at \(n=2\), supplies an actual
\(p^2\)-torsion comparison with different characters.
The manuscript index is therefore exactly two.
\end{proof}

\begin{remark}[Scalar quotients and the next arithmetic boundary]
\label{rem:B-RD71-scope}
\status{Candidate scope statement; independent audits pending.}
The new point-field argument deliberately permits scalar
residual actions after the first high quotient.
Its decisive inputs are the actual high-target inverse
valuation, its prime-to-\(p\) numerator, the common degree-\(q\)
point field, and the exact lower-break inequality.
It does not identify a high line from a scalar residual
representation alone.
The second quotient only preserves actual \(p\)-torsion;
the separate first-level lemma supplies precisely the
recovery statement needed at that level.

At the next permitted index \(e=2N=2(p^2-1)\), the small
interval can contain \(b_1=1\), \(b_2=2p-1\).
Their tangent exponents may satisfy the common-norm and
common-high-line congruences without being Frobenius conjugate,
and \(q b_1=q<e\) prevents use of
Proposition~\ref{prop:B-RD71-separated-high}.
This is an unexcluded arithmetic pattern, not a formal-group
existence assertion. No absolute index divisible by \(p\)
is treated in this subsection.
The retained exact supplement
\path{code/reduction_descent/verify_marked_quotient_sector.py}
reconstructs the inverse Newton polygons and the exact
point-field normalizations with rational arithmetic.
Across \(132\) sectors at \(p=5,7,11,13,17,19\), it checks
\(459{,}096\) first-level prime-parameter breaks,
\(29{,}376\) marked quotient polygons,
\(457{,}344\) low-line character weights, and
\(15{,}945\) separated-high point-field records, including
all six endpoint scalar pairs in these primes.
It also checks \(578{,}094\) tame mixed point-field inequalities
and the complete stated common-norm and high/low divisibility
case split. The author run passed in
\(1.3819726219953736\) seconds at 21:59:24 UTC.
The program, full output, scalar normalization samples and
pre-certificate proof are retained. No executed check failed.
The next arithmetic patterns are retained with their limited
scope; no formal-group existence conclusion is inferred.
These finite checks supplement the complete all-prime proofs.
\end{remark}
